\documentclass[10pt,a4paper]{amsart}
\usepackage[margin=19mm]{geometry}
\usepackage{amsmath,amssymb,mathtools}
\usepackage[scr=rsfs,cal=euler]{mathalfa}
\usepackage{xfrac}
\usepackage[unicode,hidelinks,bookmarks=false]{hyperref}
\newcommand{\ie}{i.e.\ }
\newcommand{\cf}{cf.\ }
\newcommand{\resp}{respectively\ }
\newcommand{\Subsection}{\subsection}
\providecommand{\nobreakdash}{\nobreak\hbox{-}}
\setlength{\emergencystretch}{2em}
\multlinegap0pt
\usepackage{bm}
\usepackage{natbib}
\usepackage{amssymb}
\usepackage{siunitx}
\usepackage{multirow}
\usepackage{float}
\usepackage[shortlabels]{enumitem}
\usepackage{caption}
\usepackage{booktabs}
\newtheorem{assumption}{Assumption}
\newtheorem{lemma}{Lemma}
\title{Nonparametric Proportional Hazards Model with Differential Regularization Applied to Spatial Survival Data}
\author{Lorenzo Tedesco, Francesco Finazzi}
\date{Source: arXiv:2410.13420v5 (CC BY 4.0)}
\begin{document}
\maketitle

\noindent\emph{Source and licence.} Nonparametric Proportional Hazards Model with Differential Regularization Applied to Spatial Survival Data. Version arXiv:2410.13420v5 (CC BY 4.0), distributed by arXiv under the Creative Commons Attribution 4.0 International licence, http://creativecommons.org/licenses/by/4.0/, as recorded in the arXiv licence metadata for this exact version and shown on the abstract page https://arxiv.org/abs/2410.13420v5. Full licence notice: \texttt{licenses/arxiv-2410.13420-CC-BY-4.0.txt}.\par\smallskip

\noindent\textit{Editorial scope.} Counted unit: the article's lemma lem:Q\_continuous (source0.tex lines 1647--1662). The article's complete original proof of it is delivered with it (source0.tex lines 1663--1724, one \texttt{proof} environment of 62 lines). No mathematics is written, repaired or re-derived: the statement and the proof are the article's own lines, in the article's own notation, and no retained line is substituted or re-flowed.  Retained uncounted as the source-local context the proof needs: lines 138--140; lines 150--153; lines 401--406.  Nothing was omitted from any retained span, so the package carries no omission record.  No retained line contains a citation, so the wrapper carries no bibliography and no bibliography entry.  No retained line contains a figure environment, an image-inclusion command or drawing code, so no image is delivered and the package declares no asset dependency.  Editorial formatting only, in addition: an AMS \texttt{amsart} preamble that restates the article's own declarations and macros used by the retained lines, namely the environments \texttt{assumption}, \texttt{lemma} on separate counters, together with the standard packages those lines need.  The retained environments print in this document as Assumption 1, Lemma 1 in order of appearance, whereas the article prints the same items with the numbering of its own full document.  The complete native source of the article as distributed in the arXiv e-print for this exact version is bundled unchanged as \texttt{sources/166985/source0.tex}.
\begin{equation}\label{eq:identification_h}
\int_{\Omega} h_0(\bm{p})\,\mathrm{d}\bm{p} = 0,
\end{equation}

\begin{assumption}\label{assumption:identifiability}
	(i) Conditional on $\bm{X}$ and $\bm{P}$, the censoring variable $C$ is independent of $T$.
	(ii) There exists $\tau > 0$ such that $\mathbb{P}(T \ge \tau) > 0$ and $\mathbb{P}(C \ge \tau) > 0$, ensuring that the support of $Y$ is non-degenerate. $(iii)$ The true coefficient vector \(\bm\beta_0\) is an interior point of a compact set \(\mathcal B \subset \mathbb{R}^b\). $(iv)$ The matrix \(\mathbb E[\bm X\bm X^\top]\) is positive definite. $(v)$ The spatial effect $h_0$ is smooth and satisfies \eqref{eq:identification_h}. %$(vi)$ The random variable $\bm{P}$ has a probability density function $f_{\bm{P}}(\mathbf{p})$ such that $0 < m_{\bm{P}} \le f_{\bm{P}}(\mathbf{p}) \le M_{\bm{P}}$ for all $\mathbf{p} \in \Omega$, where $m_{\bm{P}}$ and $M_{\bm{P}}$ are positive constants.
\end{assumption}

\begin{assumption}\label{assumption:density_P}
	(i) The support $\mathcal{X}$ of the covariates $\bm X$ is bounded, i.e.,
	\(
	\sup_{\bm x\in \mathcal{X}} \|\bm x\|\le M_{\mathcal{X}}.
	\) (ii) The triangulation $\mathcal{T}_\eta$, indexed by $\eta$, is assumed to be quasi-uniform.
\end{assumption}

	\begin{lemma}\label{lem:Q_continuous}
		Equip $\Theta=\mathcal B\times\mathcal H$ with the metric
		\(d\bigl((\bm \beta,h),(\tilde{\bm \beta},\tilde h)\bigr)=
		\|\bm \beta-\tilde{\bm\beta}\|+\|h-\tilde h\|_{\infty}\).
		Under Assumptions~\ref{assumption:identifiability}–\ref{assumption:density_P},
		the mapping
		\[
		Q(\theta)=\mathbb E\!\Bigl[
		\delta\bigl(\bm X^{\!\top}\bm \beta
		+h(\bm P)
		-\log s^{(0)}(\bm \beta,h,Y)\bigr)
		\Bigr],
		\qquad \theta=(\bm \beta,h)\in\Theta,
		\]
		is continuous on \(\Theta\).
	\end{lemma}

     \begin{proof}
    Write $s^{(0)}(\bm\beta,h,y)=\mathbb E\!\left[I(Y\ge y)\exp\{\bm X^{\!\top}\bm\beta+h(\bm P)\}\right]$.
    By Assumption~\ref{assumption:density_P}(i) there is $M_{\mathcal X}<\infty$ with
    $\|\bm X\|\le M_{\mathcal X}$ a.s., and by the definition of $\mathcal B$ we have
    $\|\bm\beta\|\le M_{\mathcal B}$. Moreover, by Sobolev embedding in $d=2$
    applied to $\mathcal H$ (bounded in $H^2$), there exists $M_\infty<\infty$ such that
    $\|h\|_\infty\le M_\infty$ for all $h\in\mathcal H$. Set
    \[
    C_*=M_{\mathcal B}M_{\mathcal X}+M_\infty.
    \]
    Then for all $(\bm\beta,h)\in\Theta$ and all $y\ge 0$,
    \begin{equation}\label{eq:s-bounds}
    e^{-C_*}\,\mathbb P(Y\ge y)\ \le\ s^{(0)}(\bm\beta,h,y)\ \le\ e^{C_*}\,\mathbb P(Y\ge y),
    \end{equation}
    since $e^{\bm X^\top\bm\beta+h(\bm P)}\in[e^{-C_*},e^{C_*}]$ a.s.  In particular,
    $s^{(0)}(\bm\beta,h,Y)>0$ a.s.
    
    Fix two parameters $\theta_1=(\bm\beta_1,h_1)$ and $\theta_2=(\bm\beta_2,h_2)$.
    By the mean–value theorem for the exponential and the bound above,
    \[
    \begin{aligned}
    \big|s^{(0)}(\theta_1,y)-s^{(0)}(\theta_2,y)\big|
    &= \left|\mathbb E\!\left[I(Y\ge y)\big(e^{\eta_1}-e^{\eta_2}\big)\right]\right| \\
    &\le \mathbb E\!\left[I(Y\ge y)\,e^{\max\{\eta_1,\eta_2\}}
    \left|(\bm X^\top(\bm\beta_1-\bm\beta_2))+(h_1-h_2)(\bm P)\right|\right] \\
    &\le e^{C_*}\,\mathbb P(Y\ge y)\,\big(M_{\mathcal X}\|\bm\beta_1-\bm\beta_2\|
    +\|h_1-h_2\|_\infty\big),
    \end{aligned}
    \]
    where $\eta_j=\bm X^\top\bm\beta_j+h_j(\bm P)$.
    Using \eqref{eq:s-bounds} and the mean–value theorem for $\log$,
    \begin{align*}
    \big|\log s^{(0)}(\theta_1,y)-\log s^{(0)}(\theta_2,y)\big|
    \ &\le\ \frac{\big|s^{(0)}(\theta_1,y)-s^{(0)}(\theta_2,y)\big|}
    {\min\{s^{(0)}(\theta_1,y),s^{(0)}(\theta_2,y)\}}\\
    \ &\le\ e^{2C_*}\,\big(M_{\mathcal X}\|\bm\beta_1-\bm\beta_2\|
    +\|h_1-h_2\|_\infty\big),
    \end{align*}
    where the factor $\mathbb P(Y\ge y)$ cancels, so the bound is uniform in $y$.
    
    Define, for $z=(Y,\delta,\bm X,\bm P)$ and $\theta=(\bm\beta,h)$,
    \[
    \zeta(z;\theta)=\delta\Big(\bm X^\top\bm\beta+h(\bm P)-\log s^{(0)}(\bm\beta,h,Y)\Big).
    \]
    Then for any $\theta_1,\theta_2\in\Theta$,
    \[
    \begin{aligned}
    \big|\zeta(z;\theta_1)-\zeta(z;\theta_2)\big|
    &\le \delta\Big(\big|\bm X^\top(\bm\beta_1-\bm\beta_2)\big|+\big|(h_1-h_2)(\bm P)\big|
    +\big|\log s^{(0)}(\theta_1,Y)-\log s^{(0)}(\theta_2,Y)\big|\Big)\\
    &\le \delta\,L_*\Big(\|\bm\beta_1-\bm\beta_2\|+\|h_1-h_2\|_\infty\Big),
    \end{aligned}
    \]
    with $L_*=(1+e^{2C_*})\max\{M_{\mathcal X},1\}$, which is deterministic and finite.
    Taking expectations and using $\delta\le 1$,
    \[
    |Q(\theta_1)-Q(\theta_2)|
    \le \mathbb E\big[\,|\zeta(z;\theta_1)-\zeta(z;\theta_2)|\,\big]
    \le L_*\,d(\theta_1,\theta_2).
    \]
    Thus $Q$ is globally Lipschitz on $(\Theta,d)$.
    \end{proof}

\end{document}
